\chapter{Triglycerides}
\section*{What Is This Test?}
Triglycerides (TG) are glycerol esters of three fatty acid molecules and are the primary form of lipid storage in adipose tissue and the primary lipid transported in very-low-density lipoprotein (VLDL) particles. Triglycerides are measured as part of the fasting lipid panel. Triglyceride-rich lipoproteins (VLDL and chylomicrons) are independently associated with increased risk of atherosclerotic cardiovascular disease, although the causal relationship is less robust than for LDL-C. Severe hypertriglyceridemia (TG >500 mg/dL, particularly >1,000 mg/dL) is a major risk factor for acute pancreatitis (likely due to hydrolysis of TG by pancreatic lipase, releasing free fatty acids that cause pancreatic microvascular injury, ischemia, and acinar cell necrosis). Causes of hypertriglyceridemia: (1) primary/genetic: familial hypertriglyceridemia, familial combined hyperlipidemia, lipoprotein lipase deficiency (familial chylomicronemia syndrome --- TG >1,000 mg/dL), and apo C-II deficiency, (2) secondary: obesity, insulin resistance and diabetes mellitus (type 2), excessive alcohol consumption, high-carbohydrate diets, physical inactivity, medications (estrogen/tamoxifen, beta-blockers, thiazide diuretics, glucocorticoids, atypical antipsychotics [clozapine, olanzapine], retinoids, cyclosporine, tacrolimus), hypothyroidism, chronic kidney disease, nephrotic syndrome, and pregnancy.
\section*{Alternative Names}
TG; Triglyceride Level; Serum Triglycerides; Fasting Triglycerides
\section*{Principle}
Triglycerides are measured by an enzymatic colorimetric method: lipase hydrolyzes triglycerides to glycerol and free fatty acids; glycerol kinase phosphorylates glycerol to glycerol-3-phosphate; glycerol-3-phosphate oxidase oxidizes to dihydroxyacetone phosphate and H2O2; peroxidase catalyzes the production of a colored product proportional to TG concentration. Fasting (9--12 hours) is required for accurate Friedewald LDL calculation and TG assessment.
\section*{History}
The recognition of triglycerides as the principal storage form of metabolic energy dates to the late 19th century. In 1967, Fredrickson, Levy, and Lees introduced a phenotyping system for the hyperlipoproteinemias that separated familial chylomicronemia (type I) from other familial hypertriglyceridemias, laying the groundwork for clinical classification. In 1972, Friedewald and colleagues published the widely used formula estimating LDL cholesterol from measured total cholesterol, HDL cholesterol, and triglycerides, which made triglyceride measurement integral to the lipid panel. The National Cholesterol Education Program Adult Treatment Panel III report of 2001 codified the current triglyceride categories (<150, 150--199, 200--499, and $\geq$500 mg/dL) that remain in clinical use. Large outcomes trials of fibrates and omega-3 fatty acids, including REDUCE-IT (2019) with icosapent ethyl, subsequently refocused attention on persistent hypertriglyceridemia as a modifiable cardiovascular risk factor.
\section*{Why This Test Is Ordered}
\begin{itemize}
  \item Routine cardiovascular risk assessment as part of the fasting lipid panel (total cholesterol, LDL, HDL, and triglycerides)
  \item Evaluation of suspected familial hyperlipidemia, including familial chylomicronemia syndrome (lipoprotein lipase deficiency) presenting with eruptive xanthomas, lipemia retinalis, and pancreatitis
  \item Assessment of acute pancreatitis risk, since triglyceride levels $>$500 mg/dL, and particularly $>$1,000 mg/dL, are an established trigger of acute pancreatitis
  \item Monitoring of response to lifestyle modification and pharmacotherapy (fibrates, prescription omega-3 fatty acids, statins)
  \item Evaluation of secondary causes of hypertriglyceridemia: poorly controlled diabetes, hypothyroidism, nephrotic syndrome, chronic kidney disease, alcohol excess, and culprit medications
  \item Component of metabolic syndrome assessment, where elevated triglycerides accompany central obesity, hypertension, low HDL, and impaired fasting glucose
\end{itemize}
\section*{Is This a Primary or Secondary Test}
Triglycerides are measured as part of the primary screening lipid panel for cardiovascular risk assessment. When severely elevated ($>$500 mg/dL), the triglyceride measurement becomes the primary diagnostic test for pancreatitis risk and for identifying familial chylomicronemia syndrome.
\section*{How This Test Is Performed}
Triglycerides are measured on a serum or plasma sample obtained after a 9--12 hour fast. A standard venipuncture sample is collected in a plain serum tube or an EDTA tube and run on an automated analyzer using the enzymatic colorimetric method described above. The assay measures total triglycerides, including chylomicrons and VLDL, without fractionation for routine use. A non-fasting or lipemic sample falsely elevates the result and invalidates the Friedewald LDL calculation.
\section*{What Preparation Is Needed}
Fasting is required for 9--12 hours before phlebotomy; only water is permitted during the fast. Alcohol should be avoided for 24 hours before testing because even modest intake raises triglycerides. The patient should follow a representative diet for 1--2 weeks before the test and avoid strenuous exercise for 24 hours before sampling.
\section*{Contraindications and Precautions}
There are no contraindications to standard venipuncture. Cautions: a non-fasting sample falsely elevates triglycerides and invalidates the Friedewald LDL estimate, so testing should be repeated fasting if the sample was not fasting. Acute illness, myocardial infarction, infection, surgery, and recent alcohol intake transiently elevate triglycerides, and elective testing should be deferred several weeks. Glycerol contamination from medication diluents and severe thrombocytosis can falsely raise the measured value, and prolonged sample storage before serum separation increases triglycerides.
\section*{Risks}
Minimal risk. Slight pain or bruising at the venipuncture site; rarely hematoma, infection, or vasovagal syncope.
\section*{What Happens After the Test}
Results are available within 1 day and are reported in mg/dL. The value is interpreted with the full lipid panel, and cardiovascular risk is estimated with pooled cohort equations. When triglycerides are very high ($\geq$500 mg/dL), strict dietary fat restriction and prompt pharmacologic therapy are initiated to reduce pancreatitis risk, and secondary causes (diabetes, hypothyroidism, alcohol, nephrotic syndrome) are investigated. Repeat testing is scheduled at 1--3 month intervals while treatment is adjusted.
\section*{Understanding Results}
\subsection*{Normal (<150 mg/dL)}
Optimal cardiovascular and metabolic health. Low risk of pancreatitis. TG <75 mg/dL in children 0--9 years.
\subsection*{Borderline (150--199 mg/dL)}
Mild hypertriglyceridemia. Lifestyle modifications: weight loss, dietary fat/carbohydrate reduction, aerobic exercise, and alcohol limitation.
\subsection*{High (200--499 mg/dL)}
Moderate hypertriglyceridemia. Lifestyle modification + pharmacotherapy (fibrates [fenofibrate, gemfibrozil], omega-3 fatty acids [icosapent ethyl], statins).
\subsection*{Very High ($\geq$500 mg/dL)}
Severe hypertriglyceridemia. Risk of acute pancreatitis. TG >1,000 mg/dL: acute pancreatitis risk, chylomicronemia syndrome (eruptive xanthomas, lipemia retinalis, pancreatitis). Requires aggressive pharmacotherapy (fibrates, high-dose prescription omega-3 fatty acids), extreme dietary fat restriction (<15--20 g/day), and avoidance of alcohol and culprit medications.
\section*{Normal Results}
Normal: <150 mg/dL. Confirm fasting sample for TG >400 (Friedewald equation invalid >400). Non-fasting TG <200 mg/dL acceptable for screening.
\section*{Follow-up Tests}
Fasting lipid panel, lipoprotein electrophoresis (chylomicrons present in TG >1,000 mg/dL), apoB, Lp(a), serum amylase/lipase (if pancreatitis suspected), TSH (hypothyroidism), HbA1c, and liver function.
\section*{References}
\begin{enumerate}
  \item Miller M, Stone NJ, Ballantyne C, et al. Triglycerides and cardiovascular disease: a scientific statement from the AHA. Circulation. 2011;123(20):2292--2333.
  \item Berglund L, Brunzell JD, Goldberg AC, et al. Evaluation and treatment of hypertriglyceridemia. J Clin Endocrinol Metab. 2012;97(9):2969--2989.
  \item Grundy SM, Stone NJ, Bailey AL, et al. 2018 AHA/\allowbreak{}ACC/\allowbreak{}AACVPR/\allowbreak{}AAPA/\allowbreak{}ABC/\allowbreak{}ACPM/\allowbreak{}ADA/\allowbreak{}AGS/\allowbreak{}APhA/\allowbreak{}ASPC/\allowbreak{}NLA/\allowbreak{}PCNA Guideline on the Management of Blood Cholesterol: a report of the American College of Cardiology/American Heart Association Task Force on Clinical Practice Guidelines. Circulation. 2019;139(25):e1082--e1143.
  \item Friedewald WT, Levy RI, Fredrickson DS. Estimation of the concentration of low-density lipoprotein cholesterol in plasma, without use of the preparative ultracentrifuge. Clin Chem. 1972;18(6):499--502.
  \item Bhatt DL, Steg PG, Miller M, et al. Cardiovascular risk reduction with icosapent ethyl for hypertriglyceridemia. N Engl J Med. 2019;380(1):11--22.
\end{enumerate}
\clearpage
